\documentclass[11pt]{article}

\usepackage[affil-it]{authblk}
\renewcommand\Authfont{\normalsize}
\renewcommand\Affilfont{\small\normalfont}
\setlength{\affilsep}{0.5em}

\usepackage[a4paper,margin=1in]{geometry}
\usepackage{amsmath,amssymb,amsthm,mathtools}
\usepackage{booktabs}
\usepackage{microtype}
\usepackage[numbers,sort&compress]{natbib}
\usepackage[colorlinks=true,linkcolor=blue,citecolor=blue,urlcolor=blue]{hyperref}
\usepackage[nameinlink,capitalise,noabbrev]{cleveref}

\newtheorem{theorem}{Theorem}
\newtheorem{proposition}{Proposition}
\newtheorem{lemma}{Lemma}
\newtheorem{corollary}{Corollary}
\newtheorem{remark}{Remark}
\newtheorem{definition}{Definition}
\newtheorem{example}{Example}

\newcommand{\R}{\mathbb R}
\DeclareMathOperator{\rank}{rank}
\DeclareMathOperator{\im}{im}
\DeclareMathOperator{\kerop}{ker}

\title{Exact Lifting of Coarse Histories\\
in Periodic Upwind Finite Volumes}

\author[1]{Antonis Polemitis}
\author[2,*]{Nicholas Christakis}
\author[2]{Dimitris Drikakis}

\affil[1]{UNIC Evolve, University of Nicosia, Nicosia CY-2417, Cyprus}
\affil[2]{Institute for Advanced Modelling and Simulation,
University of Nicosia, Nicosia CY-2417, Cyprus}
\affil[*]{Corresponding author: \texttt{christakis.n@unic.ac.cy}}

\hypersetup{
  pdftitle={Exact Lifting of Coarse Histories in Periodic Upwind Finite Volumes},
  pdfauthor={Antonis Polemitis, Nicholas Christakis, and Dimitris Drikakis},
  pdfsubject={Finite-volume coarse-history lifting, realizability, and minimal child counts},
  pdfkeywords={finite-volume methods, coarse graining, realizability, upwind advection, behavioral systems}
}

\begin{document}
\maketitle

\begin{abstract}
A prescribed coarse finite-volume history is not automatically compatible
with a fine-grid realization and a declared fine-grid update.  This short
note studies the inverse, or image-space, counterpart of finite-horizon
predictive memory.  Instead of asking which coarse history is observed from a
given fine state, we ask which parent-average histories can be lifted to a
fine-grid initial state.

For \(P\ge2\) parent cells and periodic scalar advection discretized by
first-order positive upwind finite volumes, we give a complete constructive
characterization.  For an \(L\)-step history, every mass-conserving parent
trajectory is realizable if and only if the number of fine children per
parent satisfies \(r\ge L+1\).  We construct a normalized linear right
inverse of the observation map, canonical after fixing the zero-mean
flux-gauge and zero-interior convention.  We also describe all fine-grid
lifts of every realizable short-horizon history, identifying the remaining
nonuniqueness as periodic flux-gauge freedom and unexposed interior child
values.  When \(r<L+1\), arbitrary mass-conserving histories are no longer
realizable: admissible histories satisfy an explicit lag-\(r\) parent-level
recurrence induced by the fact that \(r\) fine-grid shifts equal one
parent-level shift.

All statements are algebraic over unrestricted real-valued fine states; no
positivity, boundedness, invariant-domain, or entropy constraints are
imposed.  Exact rational certificates reproduce a normalized lift, the
long-horizon recurrence, and an explicit nonnegative mass-conserving but
nonrealizable history.  The results provide a sharp discrete \(L\)-step
minimal child-count law for the declared finite-volume hierarchy.
\end{abstract}

\noindent\textbf{Keywords:}
finite-volume methods; coarse graining; trajectory realizability; upwind
advection; behavioral systems; minimal child count.

\medskip

\noindent\textbf{Mathematics Subject Classification (2020):}
65M08; 93B07; 93B03; 93B25.

\section{Introduction}
\label{sec:introduction}

Finite-volume methods evolve cell averages through numerical fluxes at cell
interfaces \cite{LeVeque2002,EymardGallouetHerbin2000}.  In coarse/fine
coupling, adaptive mesh refinement, multirate integration, reduced modelling,
and learned discretization, one may ask whether a proposed coarse trajectory
is compatible with a declared fine-grid computation
\cite{Berger1987,BergerColella1989,McCorquodaleColella2011,
OsherSanders1983,ConstantinescuSandu2007,BarSinai2019,KochkovEtAl2021}.

The companion finite-volume erasure paper studies the forward observation
problem for a periodic scalar upwind hierarchy
\cite{polemitis2026exactfinitehorizonmemoryconditioning}.  It asks: given a
fine initial state, what parent-average history is observed, how much hidden
memory is needed to predict that history, how well conditioned is the
corresponding information, and when is it dissipated?  A second companion
note studies one-step information boundaries
\cite{PolemitisChristakisDrikakis2026OneStep}: it shows that parent averages
together with complete initial parent-interface trace pairs close one
forward-Euler transition but do not generally close a completed
second-order strong-stability-preserving Runge--Kutta, equivalently Heun,
transition, and it shows that conservative one-step closure does not
determine entropy-accounting closure.

The present note asks the image-space counterpart: given a parent-average
history, when does some fine-grid initial state realize it?  The companion
forward-analysis paper determines rank, nullspace, conditioning, and
predictive memory of the fine-state-to-history map.  This note solves the
inverse problem for the same declared hierarchy: it constructs a normalized
right inverse, describes every fiber over an admissible history, and gives a
necessary-and-sufficient description of the image.

This inverse viewpoint is closely related to behavioral systems theory,
where a dynamical system is studied through the set of trajectories that it
generates rather than through one preferred state-space representation
\cite{Willems1986ExactModelling,PoldermanWillems1998,MarkovskyDorfler2021}.
It is also related to classical observability, realization, and model
reduction \cite{CohnDee1988,HoKalman1966,Moore1981,Forney2011,
Antoulas2005,BennerGugercinWillcox2015,PeherstorferWillcox2016}.  The
finite-volume-specific contribution here is an exact lifting construction:
the recurrence is induced explicitly by \(r\) fine-grid shifts producing one
parent-level translation, and every realizable short-horizon history has a
constructive fine-grid lift with a complete fiber description.

The word lifting is also standard in equation-free and multiscale
computation, where a coarse state is mapped to one or more compatible
microscopic states before short microscopic simulation is used to advance
macroscopic observables \cite{Kevrekidis2003}.  The present use is more
specific: it is not merely state lifting at one time, but exact
finite-horizon trajectory lifting for a declared finite-volume update,
including a complete affine-fiber description.

The result is sharp for \(P\ge2\).  Let \(P\) be the number of parent cells,
\(r\) the number of fine children per parent, and \(L\) the number of
requested update steps.  Then:
\begin{itemize}
\item if \(r\ge L+1\), every parent history through time \(L\) that
conserves the sum of parent averages is realizable by an unrestricted fine
state;
\item if \(r<L+1\), mass conservation alone is insufficient, and admissible
histories satisfy an explicit lag-\(r\) recurrence.
\end{itemize}
Thus
\[
  r_{\min}=L+1
\]
is a discrete \(L\)-step minimal child-count law.  It is not a
fixed-physical-time mesh-resolution law: if parent widths are held fixed
while \(r\) changes, then the fine-cell width and, at fixed Courant number,
the time step also change.

\paragraph{Contributions.}
This note makes three contributions:
\begin{enumerate}
\item It constructs a normalized linear right inverse that lifts every
mass-conserving \(L\)-step parent history to a fine-grid initial state when
\(r\ge L+1\), after fixing the zero-mean flux-gauge and zero-interior
convention.
\item It parameterizes all such lifts and determines the exact fiber
dimension, separating periodic flux-gauge freedom from unexposed interior
child degrees of freedom.
\item It characterizes longer admissible histories exactly by a sharp
lag-\(r\) recurrence and derives \(r_{\min}=L+1\) as a consequence of the
image characterization.
\end{enumerate}

All realizability statements are algebraic statements over the complete
state space \(\R^{Pr}\).  They do not assert realizability under positivity,
maximum-principle, invariant-domain, entropy, or physical admissibility
constraints.

\paragraph{Organization.}
\Cref{sec:hierarchy} defines the periodic upwind hierarchy and its history
map.  \Cref{sec:short-horizon} proves the short-horizon image-and-fiber
theorem and gives an explicit canonical lift.  \Cref{sec:long-horizon}
characterizes longer realizable histories by recurrence and proves the sharp
minimal child-count law.  \Cref{sec:certificates} records exact rational
certificate examples.  \Cref{sec:discussion} discusses scope and limitations,
and \cref{sec:conclusion} concludes.

\section{Periodic Upwind Hierarchy and History Map}
\label{sec:hierarchy}

Consider scalar advection
\[
  u_t+a u_x=0,
  \qquad a>0,
\]
on a periodic one-dimensional domain.  The first-order positive upwind
finite-volume update is
\[
  u_i^{n+1}
  =
  (1-\lambda)u_i^n+\lambda u_{i-1}^n,
  \qquad 0<\lambda\le1,
\]
with periodic indexing, where
\[
  \lambda=\frac{a\Delta t}{h}
\]
is the Courant number.

Let \(P\ge2\) be the number of parent cells, each containing \(r\ge1\) fine
children.  Then
\[
  N=Pr
\]
is the number of fine cells.  Write
\[
  x=(x_{K,j}),
  \qquad
  K=0,\ldots,P-1,
  \qquad
  j=0,\ldots,r-1,
\]
with parent indices interpreted modulo \(P\).

The parent-average restriction is
\[
  (Rx)_K=\frac1r\sum_{j=0}^{r-1}x_{K,j}.
\]
Let \(S\) be the downstream cyclic shift,
\[
  (Sx)_{K,j}
  =
  \begin{cases}
  x_{K,j-1}, & j\ge1,\\
  x_{K-1,r-1}, & j=0,
  \end{cases}
\]
and set
\[
  A_\lambda=(1-\lambda)I+\lambda S.
\]

For a horizon \(L\ge0\), define the history map
\[
  {\cal O}_{L,r}
  =
  \begin{bmatrix}
  R\\
  RA_\lambda\\
  \vdots\\
  RA_\lambda^L
  \end{bmatrix}
  :
  \R^{Pr}\to(\R^P)^{L+1}.
\]
Hence
\[
  {\cal O}_{L,r}x
  =
  (y^0,y^1,\ldots,y^L),
  \qquad
  y^n=RA_\lambda^n x.
\]

The dependence of \({\cal O}_{L,r}\) and of the lifting maps below on the
fixed Courant number \(\lambda\) is suppressed in the notation.

Let \(\Pi_P:\R^P\to\R^P\) denote the parent cyclic shift,
\[
  (\Pi_Py)_K=y_{K-1},
\]
and define the periodic difference operator
\[
  B=I-\Pi_P,
  \qquad
  (B\Phi)_K=\Phi_K-\Phi_{K-1}.
\]
Let
\[
  \R_0^P
  =
  \left\{
  z\in\R^P:\mathbf 1_P^\top z=0
  \right\}.
\]

The operator \(B\) maps \(\R^P\) into \(\R_0^P\), and its kernel on
\(\R^P\) is the span of \(\mathbf 1_P\).  Its restriction to the zero-mean
subspace,
\[
  B_0:=B|_{\R_0^P}:\R_0^P\to\R_0^P,
\]
is invertible.  We denote the inverse by \(B_0^{-1}\).  Thus
\(B_0^{-1}g\) is the unique zero-mean register \(\Phi\) satisfying
\[
  B\Phi=g
\]
for a given \(g\in\R_0^P\).  Indeed, \(B\Phi=0\) implies that \(\Phi\) is
constant, and the only constant vector in \(\R_0^P\) is zero.

A concrete formula is obtained as follows.  Set
\[
  \widehat\Phi_0=0,
  \qquad
  \widehat\Phi_K=\sum_{m=1}^{K}g_m,
  \qquad K=1,\ldots,P-1.
\]
Then the zero-mean solution is
\[
  B_0^{-1}g
  =
  \widehat\Phi
  -
  \frac1P\left(\sum_{K=0}^{P-1}\widehat\Phi_K\right)\mathbf 1_P.
\]

\begin{definition}[Mass-conserving parent histories]
For \(L\ge0\), define
\[
  {\cal M}_{P,L}
  =
  \left\{
  (y^0,\ldots,y^L)\in(\R^P)^{L+1}:
  \mathbf 1_P^\top y^n=\mathbf 1_P^\top y^0
  \text{ for }n=1,\ldots,L
  \right\}.
\]
\end{definition}

Every history in \(\im{\cal O}_{L,r}\) lies in \({\cal M}_{P,L}\), since the
upwind update conserves the fine-grid sum.

\section{Short-Horizon Image, Fibers, and Canonical Lifting}
\label{sec:short-horizon}

Assume
\[
  0\le L\le r-1.
\]
For \(L\ge1\), define right-collar variables
\[
  c_{K,j}=x_{K,r-1-j},
  \qquad
  j=0,\ldots,L-1.
\]
The outgoing upwind flux through the right face of parent \(K\) at time \(t\)
is
\[
  \Phi_{K,t}
  =
  \lambda(A_\lambda^t x)_{K,r-1}
  =
  \lambda
  \sum_{j=0}^{t}
  \binom{t}{j}(1-\lambda)^{t-j}\lambda^j c_{K,j},
  \qquad
  t=0,\ldots,L-1.
\]
Write this map as
\[
  \begin{bmatrix}
  \Phi_{K,0}\\
  \Phi_{K,1}\\
  \vdots\\
  \Phi_{K,L-1}
  \end{bmatrix}
  =
  T_{L,\lambda}
  \begin{bmatrix}
  c_{K,0}\\
  c_{K,1}\\
  \vdots\\
  c_{K,L-1}
  \end{bmatrix}.
\]
The matrix \(T_{L,\lambda}\) is lower triangular with diagonal entries
\[
  \lambda,\lambda^2,\ldots,\lambda^L,
\]
and hence is invertible for \(0<\lambda\le1\).

The inverse is explicit.  If
\[
  \phi_K=(\Phi_{K,0},\ldots,\Phi_{K,L-1})^\top,
  \qquad
  c_K=T_{L,\lambda}^{-1}\phi_K,
\]
then for \(t=0,\ldots,L-1\),
\[
  c_{K,t}
  =
  \lambda^{-(t+1)}
  \sum_{j=0}^{t}
  \binom{t}{j}
  \bigl[-(1-\lambda)\bigr]^{t-j}
  \Phi_{K,j},
  \qquad
  t=0,\ldots,L-1.
\]
This is the binomial inverse of the lower-triangular collar-to-flux map.

\begin{theorem}[Short-horizon image, fibers, and canonical lift]
\label{thm:short-realization}
Let \(P\ge2\), \(r\ge1\), \(0<\lambda\le1\), and \(0\le L\le r-1\).  Over
the unrestricted fine state space \(\R^{Pr}\),
\[
  \im{\cal O}_{L,r}={\cal M}_{P,L}.
\]
Moreover:
\begin{enumerate}
\item There exists a normalized linear right inverse, canonical relative to the
zero-mean flux-gauge and zero-interior convention,
\[
  {\cal L}_{L,r}:{\cal M}_{P,L}\to\R^{Pr}
\]
such that
\[
  {\cal O}_{L,r}{\cal L}_{L,r}=I_{{\cal M}_{P,L}}.
\]

\item For every \(Y\in{\cal M}_{P,L}\), the fiber
\[
  {\cal O}_{L,r}^{-1}(Y)
\]
is a nonempty affine space of dimension
\[
  P(r-L-1)+L.
\]

\item The fiber degrees of freedom consist of \(L\) periodic flux-gauge
constants, one for each transition, and
\[
  P(r-L-1)
\]
unexposed interior child values.
\end{enumerate}
\end{theorem}

\begin{proof}
For \(L=0\), the statement is immediate: every
\[
  y^0\in\R^P
\]
is lifted canonically by setting
\[
  x_{K,0}=r y_K^0,
  \qquad
  x_{K,j}=0\quad(j=1,\ldots,r-1).
\]
The fiber dimension is \(P(r-1)\), as stated.

Now let \(L\ge1\).  The inclusion
\[
  \im{\cal O}_{L,r}\subseteq{\cal M}_{P,L}
\]
follows from conservation.

Conversely, let
\[
  Y=(y^0,\ldots,y^L)\in{\cal M}_{P,L}.
\]
For each transition \(t=0,\ldots,L-1\), define
\[
  g^t=r(y^t-y^{t+1}).
\]
Since \(Y\) conserves the sum of parent averages,
\[
  \mathbf 1_P^\top g^t=0.
\]
Choose the unique zero-mean flux register
\[
  \Phi^t=B_0^{-1}g^t\in\R_0^P.
\]
Then
\[
  B\Phi^t=r(y^t-y^{t+1}),
\]
or equivalently
\[
  y^{t+1}=y^t-\frac1rB\Phi^t.
\]

For each parent \(K\), define the flux history vector
\[
  \phi_K
  =
  (\Phi_{K,0},\ldots,\Phi_{K,L-1})^\top.
\]
Set
\[
  c_K=T_{L,\lambda}^{-1}\phi_K,
\]
and assign
\[
  x_{K,r-1-j}=c_{K,j},
  \qquad j=0,\ldots,L-1.
\]
Assign all unexposed interior children
\[
  x_{K,j}=0,
  \qquad j=1,\ldots,r-L-1,
\]
and finally set
\[
  x_{K,0}
  =
  r y_K^0-\sum_{j=1}^{r-1}x_{K,j}.
\]
This defines the normalized lift \({\cal L}_{L,r}Y\).  Its parent averages
equal \(y^0\), and its outgoing fluxes equal \(\Phi^t\) for
\(t=0,\ldots,L-1\).  The parent-update identity above then gives
\[
  {\cal O}_{L,r}{\cal L}_{L,r}Y=Y.
\]

It remains to characterize the fibers.  Let \(d\in\ker{\cal O}_{L,r}\).
Then
\[
  Rd=0
\]
and the parent update shows that each outgoing-flux vector associated with
\(d\) must satisfy
\[
  B\Phi^t(d)=0,
  \qquad t=0,\ldots,L-1.
\]
Hence
\[
  \Phi^t(d)=\gamma_t\mathbf 1_P
\]
for some independent scalar gauge constants
\[
  \gamma_0,\ldots,\gamma_{L-1}.
\]
Inverting \(T_{L,\lambda}\), the corresponding collar perturbations are
identical in every parent and are uniquely determined by
\[
  T_{L,\lambda}^{-1}
  (\gamma_0,\ldots,\gamma_{L-1})^\top.
\]
The unexposed interior children
\[
  d_{K,j},
  \qquad
  K=0,\ldots,P-1,
  \quad
  j=1,\ldots,r-L-1,
\]
are arbitrary.  The values \(d_{K,0}\) are then uniquely fixed by the
conditions \(Rd=0\).  Thus the kernel is parameterized by \(L\) gauge
constants and \(P(r-L-1)\) free interior values, giving dimension
\[
  L+P(r-L-1).
\]
Every fiber is therefore the translate
\[
  {\cal O}_{L,r}^{-1}(Y)
  =
  {\cal L}_{L,r}Y+\ker{\cal O}_{L,r}.
\]
\end{proof}

\begin{example}[Normalized short-horizon lift]
\label{ex:canonical-lift}
Let
\[
  P=2,
  \qquad
  L=1,
  \qquad
  r=2,
  \qquad
  \lambda=\frac12.
\]
Consider the mass-conserving one-step parent history
\[
  y^0=(0,0),
  \qquad
  y^1=(1,-1).
\]
Then
\[
  g^0=r(y^0-y^1)=(-2,2).
\]
The unique zero-mean interface register satisfying
\[
  B\Phi^0=g^0
\]
is
\[
  \Phi^0=(-1,1).
\]
Since
\[
  \Phi_{K,0}=\frac12c_{K,0},
\]
the right-collar values are
\[
  c_{0,0}=-2,
  \qquad
  c_{1,0}=2.
\]
The normalized lift is
\[
  x=(2,-2\mid -2,2).
\]
Direct calculation gives
\[
  Rx=(0,0),
  \qquad
  RA_{1/2}x=(1,-1).
\]
Thus
\[
  {\cal O}_{1,2}x=(y^0,y^1).
\]
\end{example}

At the sharp child count
\[
  r=L+1,
\]
the fiber dimension in \cref{thm:short-realization} reduces to
\[
  L.
\]
Thus the only nonuniqueness in a short-horizon lift is the \(L\)-dimensional
periodic flux gauge.

\begin{remark}[Algebraic lifting versus stable lifting]
The theorem establishes exact algebraic surjectivity.  It does not assert
uniform boundedness or numerical stability of the right inverse.  The explicit
binomial-inverse formula above contains inverse powers of \(\lambda\), so
prescribed \(O(1)\) changes in parent histories may require large collar
values when \(\lambda\) is small. The inverse
\(B_0^{-1}\) also reflects the conditioning of the periodic difference
operator and depends on \(P\).  Thus exact realizability and stable or
physically moderate realizability are distinct questions.
\end{remark}

\section{Long-Horizon Realizability}
\label{sec:long-horizon}

When \(L\ge r\), mass conservation is no longer the only constraint.  The
parent history must satisfy a linear recurrence induced by the identity
\[
  \bigl(A_\lambda-(1-\lambda)I\bigr)^r=\lambda^rS^r.
\]
Since
\[
  RS^r=\Pi_PR,
\]
this identity closes at the parent level after \(r\) fine-grid shifts.

Define the all-horizon invisible subspace
\[
  {\cal N}_r
  =
  \left\{
  x\in\R^{Pr}:
  x_{K,j}=c_j\ \text{for all }K,\quad
  \sum_{j=0}^{r-1}c_j=0
  \right\}.
\]
It has dimension \(r-1\).

\begin{theorem}[Long-horizon image and fibers]
\label{thm:long-realization}
Let \(P\ge2\), \(r\ge1\), \(0<\lambda\le1\), and \(L\ge r-1\).  A parent
history
\[
  Y=(y^0,\ldots,y^L)\in(\R^P)^{L+1}
\]
belongs to \(\im{\cal O}_{L,r}\) if and only if:
\begin{enumerate}
\item the first \(r\) parent states have equal sums,
\[
  \mathbf 1_P^\top y^0
  =
  \cdots
  =
  \mathbf 1_P^\top y^{r-1};
\]
\item if \(L\ge r\), then for every \(n=0,\ldots,L-r\),
\begin{equation}
  y^{n+r}
  =
  \lambda^r\Pi_P y^n
  -
  \sum_{j=0}^{r-1}
  \binom{r}{j}
  \bigl[-(1-\lambda)\bigr]^{r-j}
  y^{n+j}.
  \label{eq:history-recurrence}
\end{equation}
\end{enumerate}
For every realizable history,
\[
  {\cal O}_{L,r}^{-1}(Y)
  =
  x_\star+{\cal N}_r
\]
for any one lift \(x_\star\), and hence
\[
  \dim{\cal O}_{L,r}^{-1}(Y)=r-1.
\]
\end{theorem}

\begin{proof}
Suppose first that
\[
  Y={\cal O}_{L,r}x.
\]
Conservation gives the equality of parent sums.  Expanding
\[
  \bigl(A_\lambda-(1-\lambda)I\bigr)^r=\lambda^rS^r
\]
gives
\[
  \sum_{j=0}^{r}
  \binom{r}{j}
  \bigl[-(1-\lambda)\bigr]^{r-j}
  A_\lambda^j
  =
  \lambda^rS^r.
\]
Multiplying by \(RA_\lambda^n\) and using
\[
  RS^r=\Pi_PR
\]
gives
\[
  \sum_{j=0}^{r}
  \binom{r}{j}
  \bigl[-(1-\lambda)\bigr]^{r-j}
  y^{n+j}
  =
  \lambda^r\Pi_Py^n.
\]
Solving for \(y^{n+r}\) gives \eqref{eq:history-recurrence}.

Conversely, suppose the stated sum condition and recurrence hold.  By
\cref{thm:short-realization}, applied at horizon \(r-1\), there exists a
fine state \(x_\star\in\R^{Pr}\) whose first \(r\) parent states are
\[
  y^0,\ldots,y^{r-1}.
\]
The parent history generated by \(x_\star\) satisfies
\eqref{eq:history-recurrence}.  Since it agrees with \(Y\) for the first
\(r\) states, induction gives agreement through time \(L\).

It remains to identify the fiber.  Because the temporal transformation from
\[
  (R,RS,\ldots,RS^{r-1})
\]
to
\[
  (R,RA_\lambda,\ldots,RA_\lambda^{r-1})
\]
is triangular with nonzero diagonal entries, a vector
\[
  d\in\ker{\cal O}_{r-1,r}
\]
satisfies
\[
  RS^td=0,
  \qquad t=0,\ldots,r-1.
\]
For \(t=0,\ldots,r-2\),
\[
  r(RS^{t+1}d-RS^td)_K
  =
  d_{K-1,r-1-t}-d_{K,r-1-t}.
\]
Thus every child layer
\[
  j=1,\ldots,r-1
\]
is constant across parents.  Since \(Rd=0\), child layer \(j=0\) is also
constant across parents and the common child profile has zero sum.  Hence
\[
  \ker{\cal O}_{r-1,r}={\cal N}_r.
\]
The recurrence shows that no later observation reduces this kernel, so
\[
  \ker{\cal O}_{L,r}={\cal N}_r
  \qquad
  (L\ge r-1).
\]
Therefore every realizable fiber is
\[
  x_\star+{\cal N}_r
\]
and has dimension \(r-1\).
\end{proof}

\begin{corollary}[Sharp discrete minimal child-count law]
\label{cor:minimal-r}
Let \(P\ge2\) and \(L\ge0\).  For a fixed number of update steps \(L\), the
minimal number of children per parent needed to realize every
mass-conserving history in \({\cal M}_{P,L}\), over unrestricted real-valued
fine states and for \(0<\lambda\le1\), is
\[
  r_{\min}=L+1.
\]
\end{corollary}

\begin{proof}
By \cref{thm:short-realization}, \(r=L+1\) is sufficient.

Conversely, suppose \(r\le L\).  Choose a nonzero vector
\[
  v\in\R_0^P.
\]
Define a mass-conserving history through time \(L\) by
\[
  y^0=\cdots=y^{r-1}=0,
  \qquad
  y^r=v,
\]
and choose any zero-sum states for later times, for example
\[
  y^{r+1}=\cdots=y^L=0.
\]
The recurrence \eqref{eq:history-recurrence} at \(n=0\) predicts
\[
  y^r=0,
\]
because
\[
  y^0=\cdots=y^{r-1}=0.
\]
This contradicts \(y^r=v\ne0\).  Hence the history is mass-conserving but
not realizable.  Therefore \(r\le L\) is insufficient.
\end{proof}

\begin{example}[Nonnegative mass-conserving but nonrealizable history]
\label{ex:nonrealizable}
Let
\[
  P=2,
  \qquad
  r=2,
  \qquad
  L=2,
  \qquad
  \lambda=\frac12.
\]
Consider
\[
  y^0=(1,1),
  \qquad
  y^1=(1,1),
  \qquad
  y^2=(2,0).
\]
The history is nonnegative and conserves the sum of parent averages.  However,
the lag-two recurrence \eqref{eq:history-recurrence} with
\[
  y^0=y^1=(1,1)
\]
predicts
\[
  y^2=(1,1).
\]
The recurrence residual is therefore
\[
  (1,-1),
\]
so no fine initial state with two children per parent realizes this
two-step parent history.
\end{example}

The minimal child-count law is a discrete-step law.  If parent width is fixed
and \(r\) changes, then fine-cell width changes.  At fixed Courant number,
the fine time step changes as well.  Thus the law does not claim a
fixed-physical-time mesh-resolution criterion.

\section{Exact Computational Certificates}
\label{sec:certificates}

The accompanying script verifies representative finite instances of the
theorems and generates the reported exact rational examples.  The general
results are established analytically above.  All certificate calculations use
Python's \texttt{fractions.Fraction}, so no floating-point roundoff enters
the reported values.

\begin{table}[htbp]
\centering
\caption{Exact rational certificates generated by the accompanying script.
All checks are performed with Python's \texttt{fractions.Fraction}, so the
results do not depend on floating-point roundoff.  The first row constructs
the normalized short-horizon lift in \cref{ex:canonical-lift}.  The second
row checks that a generated parent history satisfies the order-\(r\)
recurrence.  The third row gives the nonnegative nonrealizable
mass-conserving history in \cref{ex:nonrealizable}.}
\label{tab:certificates}
\begin{tabular}{@{}lll@{}}
\toprule
Certificate & Parameters & Result \\
\midrule
Normalized short-horizon lift & \(P=2,\ L=1,\ r=2,\ \lambda=1/2\) & exact lift \\
Long-horizon recurrence & \(P=3,\ r=3,\ L=6,\ \lambda=1/2\) & residual zero \\
Nonrealizable nonnegative history & \(P=2,\ r=2,\ L=2,\ \lambda=1/2\) & residual \((1,-1)\) \\
\bottomrule
\end{tabular}
\end{table}

\Cref{tab:certificates} provides reproducibility checks for the finite
examples.  It is not used as a substitute for the analytical proofs.

\section{Discussion}
\label{sec:discussion}

The present note is the image-space counterpart of the finite-horizon
observation analysis in
\cite{polemitis2026exactfinitehorizonmemoryconditioning}, which
characterizes rank, nullspace, conditioning, and predictive memory of the
forward history map.  It is also complementary to the one-step
information-boundary analysis in
\cite{PolemitisChristakisDrikakis2026OneStep}.  Here we solve the inverse
problem: we characterize the image of the history map, construct a normalized
right inverse, and parameterize the fine-state fibers over every realizable
parent history.

The finite-volume-specific content is not merely that finite-dimensional
linear trajectories satisfy a recurrence.  The recurrence here has an
explicit form induced by
\[
  S^r=\text{one parent-level shift},
\]
begins sharply after the first \(r\) parent states, and is paired with a
constructive fine-grid lifting map.  The short-horizon fibers separate
periodic flux-gauge freedom from interior child values that have not yet
reached a parent interface.

This is not a validation theorem for nonlinear CFD.  It is an exact algebraic
statement for a declared periodic scalar upwind hierarchy over unrestricted
real-valued fine states.  Positivity, maximum principles, entropy
admissibility, physical state constraints, nonlinear systems, nonperiodic
boundaries, high-order reconstruction, and multistage integration require
separate analysis.

\section{Conclusion}
\label{sec:conclusion}

For \(P\ge2\), \(0<\lambda\le1\), and unrestricted real-valued fine states,
every mass-conserving parent history over \(L\) update steps in the periodic
scalar upwind hierarchy is exactly liftable to a fine-grid initial state if
and only if the number of children per parent satisfies
\[
  r\ge L+1.
\]
The lift can be chosen in a normalized canonical form by fixing the
zero-mean periodic flux gauge, inverting the collar-to-flux map, setting
unexposed interior children to zero, and using the \(j=0\) child to enforce the initial parent averages.

When \(r<L+1\), mass conservation is insufficient.  The parent history must
satisfy the explicit lag-\(r\) recurrence \eqref{eq:history-recurrence}.
The full fiber structure is also determined: short-horizon lift ambiguity
consists of periodic flux gauges and unexposed interior children, whereas
long-horizon lift ambiguity reduces to the repeated zero-mean parent profile.

The result provides a sharp discrete minimal child-count law and a
constructive image-and-fiber characterization for the periodic upwind
hierarchy. Together with the finite-horizon memory paper
\cite{polemitis2026exactfinitehorizonmemoryconditioning} and the one-step
information-boundary note \cite{PolemitisChristakisDrikakis2026OneStep},
this result provides a three-part algebraic baseline for the scalar periodic
upwind hierarchy: recursive prediction requires exact memory, one-step
closure requires the correct interface and entropy-side information, and
coarse histories themselves require sufficient subcell resolution to be
realizable.

\section*{Funding}

No external funding was received for this work.

\section*{Declaration of competing interest}

The authors declare that they have no known competing financial interests or
personal relationships that could have appeared to influence the work
reported in this paper.

\section*{CRediT authorship contribution statement}

Antonis Polemitis: Conceptualization, Methodology, Verification, Writing--review and editing.  Nicholas Christakis:
Conceptualization, Methodology, Formal analysis, Software, Verification,
Writing--original draft, Writing--review and editing.  Dimitris Drikakis:
Conceptualization, Methodology, Verification, Writing--review and editing.

\section*{Data and code availability}

The source code, exact rational certificate script, generated CSV files, and
LaTeX table fragments supporting this study are available at
\url{https://github.com/UniversityOfNicosia/certified-simulation}, under
\texttt{papers/coarse-history-realizability}.  The finite certificate checks
use exact rational arithmetic via Python's \texttt{fractions.Fraction};
hence the reported finite identities, recurrence residuals, and table entries
do not depend on floating-point roundoff.

\bibliographystyle{unsrt}
\bibliography{references}

\end{document}